\begin{definition}[Дефицит]
\(\dfc(G)=n-2\alpha'(G)\) — число вершин, не покрытых максимальным
паросочетанием. Эквивалентно \(\alpha'(G)=\dfrac{n-\dfc(G)}{2}\).
\(\dfc(G)=0\) тогда и только тогда, когда есть совершенное паросочетание.
\end{definition}
\begin{theorem}[Формула Бержа, 1958]
\[
\dfc(G)=\max_{S\subseteq V}\bigl(\odd(G-S)-|S|\bigr),\qquad
\alpha'(G)=\frac12\Bigl(n-\max_{S\subseteq V}\bigl(\odd(G-S)-|S|\bigr)\Bigr).
\]
\end{theorem}
При \(S=\varnothing\) выражение равно \(\odd(G)\ge0\), так что максимум неотрицателен.
Теорема Татта — частный случай: \(\dfc(G)=0\iff\odd(G-S)\le|S|\) для всех \(S\).
\begin{proof}
\(\ge\). Пусть \(M\) — максимальное паросочетание, \(S\subseteq V\),
\(Q_1,\ldots,Q_t\) — нечётные компоненты \(G-S\), \(t=\odd(G-S)\). В каждой \(Q_i\)
внутренние рёбра \(M\) не покрывают хотя бы одну вершину \(u_i\); она либо свободна,
либо сопоставлена вершине из \(S\). Сопоставленных не больше \(|S|\), значит,
свободных среди \(u_1,\ldots,u_t\) не меньше \(t-|S|\): \(\dfc(G)\ge\odd(G-S)-|S|\).

\(\le\). Пусть \(k=\max_S(\odd(G-S)-|S|)\). При \(k=0\) по теореме Татта есть
совершенное паросочетание и \(\dfc(G)=0\). Пусть \(k>0\).

\emph{Чётность.} Для любого \(S\): \(n-|S|=\sum|V(Q)|\) по компонентам \(G-S\),
и по модулю 2 это \(\odd(G-S)\). Значит, \(\odd(G-S)-|S|\equiv n\pmod2\),
в частности \(k\equiv n\pmod 2\).

Построим \(H\): добавим к \(G\) клику \(W\) из \(k\) новых вершин и соединим каждую
вершину \(W\) со всеми вершинами \(G\). Тогда \(v(H)=n+k\) чётно. Проверим условие
Татта для \(H\) и \(T\subseteq V(H)\):
\begin{enumerate}\tightlist
\item \(T=\varnothing\): \(H\) связен и чётен, \(\odd(H)=0\).
\item \(T\ne\varnothing\), \(W\not\subseteq T\): в \(H-T\) осталась вершина \(W\), смежная
  со всеми, поэтому \(H-T\) связен и \(\odd(H-T)\le1\le|T|\).
\item \(W\subseteq T\), \(T=W\cup S\), \(S\subseteq V\): \(H-T=G-S\) и
  \(\odd(H-T)=\odd(G-S)\le k+|S|=|T|\) по определению \(k\).
\end{enumerate}
По теореме Татта в \(H\) есть совершенное паросочетание \(N\), \(|N|=(n+k)/2\).
Выбросим из \(N\) рёбра, инцидентные \(W\) — их не больше \(k\). Останется
паросочетание \(G\) размера \(\ge\frac{n+k}{2}-k=\frac{n-k}{2}\), т.\,е.
\(\dfc(G)\le n-(n-k)=k\).
\end{proof}
\begin{graphillustration}{Граф $H$ из доказательства (здесь $n=6$, $k=2$): новые вершины $W$ образуют клику и смежны со всеми вершинами $G$. Совершенное паросочетание $H$ использует не более $k$ рёбер при $W$; остальные его рёбра — паросочетание $G$.}{fig:05-construction}
\begin{tikzpicture}[font=\small]
\node[draw,ellipse,minimum width=4.2cm,minimum height=2.2cm] (G) at (0,0) {};
\node[gi note] at (0,-1.4) {$G$};
\foreach \i/\x/\y in {1/-1.3/0.4,2/-0.3/0.6,3/0.9/0.5,4/1.2/-0.4,5/-1.1/-0.5,6/0.0/-0.5}
  \node[gi small] (g\i) at (\x,\y) {};
\node[gi vertex] (w1) at (4.2,0.8) {$w_1$};
\node[gi vertex] (w2) at (4.2,-0.8) {$w_2$};
\draw[gi edge] (w1) -- (w2);
\node[gi group,fit=(w1)(w2),label=right:{$W$, $|W|=k$}] {};
\foreach \i in {1,...,6} {\draw[gi edge,black!30] (w1) -- (g\i); \draw[gi edge,black!30] (w2) -- (g\i);}
\draw[gi match] (w1) -- (g3); \draw[gi match] (w2) -- (g4);
\draw[gi match] (g1) -- (g2); \draw[gi match] (g5) -- (g6);
\end{tikzpicture}
\end{graphillustration}


\begin{example}
Для звезды \(K_{1,3}\) с центром \(c\): \(S=\{c\}\) даёт \(\odd(G-S)-|S|=3-1=2\),
поэтому \(\dfc=2\), \(\alpha'=(4-2)/2=1\). Для графа на рис.~\ref{fig:04-tutte-set}:
\(n=6\), \(\dfc\ge3-1=2\), и действительно \(\alpha'=2\).
\end{example}
\begin{tcolorbox}[colback=cbGrayBack,colframe=cbGray!55]
\textbf{Как пользоваться.} Чтобы доказать \(\alpha'(G)\le r\), достаточно предъявить
\(S\) с \(\odd(G-S)-|S|\ge n-2r\) («свидетель»); чтобы доказать \(\alpha'(G)\ge r\) —
предъявить паросочетание. Формула гарантирует, что свидетель всегда существует.
\end{tcolorbox}
